%!TEX root=main.tex
Recent work has shown that convolutional networks can be substantially deeper, more accurate, and efficient to train if they contain shorter connections between layers close to the input and those close to the output.
In this paper, we embrace this observation and introduce the \methodnamecap{} (\methodnameshort{}), which connects each layer to every other layer in a feed-forward fashion. Whereas traditional convolutional networks with $L$ layers have $L$ connections---one between each layer and its subsequent layer---our network has $\frac{L(L+1)}{2}$ direct connections.
For each layer, the feature-maps of all preceding layers are used as inputs, and its own feature-maps are used as inputs into all subsequent layers.
\methodnameshorts{} have several compelling advantages:
they alleviate the vanishing-gradient problem, strengthen feature propagation, encourage feature reuse, and substantially reduce the number of parameters.
We evaluate our proposed architecture on four highly competitive object recognition benchmark tasks (CIFAR-10, CIFAR-100, SVHN, and ImageNet). \methodnameshorts{} obtain significant improvements over the state-of-the-art on most of them, whilst requiring less computation to achieve high performance. Code and pre-trained models are available at \url{https://github.com/liuzhuang13/DenseNet}.

%Similar to the popular ResNet architecture, \methodnameshort{} is straight-forward to optimize even with hundreds of layers---however it tends to generalize better and requires fewer parameters.


%Our  is in contrast to the traditional convolutional network layout, where preceding feature maps pass through non-linearities, or
%ResNets, where they are part of a single cumulative input.
%In contrast to the common convolutional network architecture, where each layer transforms only the output of its direct predecessor,
%\methodnameshort{} adopts a cascade structure in which feature maps extracted by preceding layers are directly used as inputs of later convolutional layers.


%We show that \methodnameshort{} with hundreds of layers can be easily trained without any optimization difficulties, like the recently proposed Residual Networks (ResNets) \cite{resnet, identity-mappings}. To achieve the same level of accuracy, \methodnameshorts{} generally requires significantly fewer parameters than ResNet. State-of-the-art results have been obtained on benchmark datasets CIFAR-10, CIFAR100 and SVHN.

%new network architecture
%allows training of networks with hundreds of layers. The features are learned in a cascade manner: in each layer, the network only learns a small
%amount of feature maps, and then pass them to next layer together with its input feature maps. Several difficulites in training very deep networks
% are solved naturally using this architecture.  Experiments on benchmark datasets (CIFAR-10, CIFAR100, SVHN) demonstrate that \methodnameshort{}
% can achieve state-of-the-art or comparable performances, while using much fewer parameters than ResNet \cite{resnet, identity-mappings}.
% %The effectiveness of \methodnameshort{} shows that the identity connections in the state-of-the-art ResNet \cite{resnet,identity-mappings} are crucial to its success, while the concrete formulation as residual learning may be incidental.In addition, \methodnameshort{} has several advantages such as implicit deep supervision and test-time memory saving. Analysis on the behavior of \methodnameshort{} helps us understand some intriguing properties of very deep networks.
